\section{Sequence-level Zeckendorf spectra and falsifiable predictions}
\label{sec:sequence_spectra}

The codon-level map $c\mapsto (V_{\mu^\ast}(c),\Delta_{\mu^\ast}(c))$ induces a natural spectrum representation for sequences.

\subsection{Definitions: \texorpdfstring{$Z$}{Z}-spectrum and uplift trace}
\label{sec:z_u_def}

Let $S$ be an RNA (or DNA) sequence and fix a reading frame.
Split $S$ into codons $(c_1,c_2,\dots,c_L)$.
Define:
\begin{itemize}
\item the \emph{$Z$-spectrum} (visible Zeckendorf level)
\[
Z(i):=V_{\mu^\ast}(c_i)\in\{0,\dots,20\};
\]
\item the \emph{uplift trace} (window-external load)
\[
U(i):=\Delta_{\mu^\ast}(c_i)=N_{\mu^\ast}(c_i)-V_{\mu^\ast}(c_i).
\]
\end{itemize}

Intuitively, $Z(i)$ measures the locally visible folding complexity in the Zeckendorf window, while $U(i)$ measures the hidden external component that is erased by folding.

\subsection{Predictions and tests}
\label{sec:predictions}

We list several directly testable predictions that do not rely on any parameter fitting beyond the unique encoding $\mu^\ast$.

\paragraph{P1 (start/stop boundary split).}
In the canonical reading frame of annotated ORFs, the start codon $\mathrm{AUG}$ and the stop codon $\mathrm{UAA}$ should produce the same boundary Zeckendorf level $Z=14$ but different uplift values:
\[
(Z,U)(\mathrm{AUG})=(14,0),\qquad (Z,U)(\mathrm{UAA})=(14,34).
\]
Quantitative variants of this check include start/stop $(Z,U)$ histograms and boundary-hit rates $B(i)$ at annotated control sites (Appendix~\ref{app:stat_tests}).

\paragraph{P2 (UGA extreme uplift).}
The stop codon $\mathrm{UGA}$ has a large uplift $U=55$ in the minimal window-$6$ model.
If context-dependent recoding (e.g.\ selenocysteine insertion at $\mathrm{UGA}$ or pyrrolysine insertion at $\mathrm{UAG}$) is sensitive to hidden loads, then the neighborhood statistics of $U(i)$ around $\mathrm{UGA}$ and $\mathrm{UAG}$ sites should differ systematically from those around $\mathrm{UAA}$ sites \cite{HatfieldGladyshev2002,SrinivasanJamesKrzycki2002,McCaughanBrownDalphinBerryTate1995}.

\paragraph{P3 (nonstandard start/stop codes).}
In organisms with nonstandard start/stop assignments, our framework predicts either:
\begin{itemize}
\item a reassignment that moves control semantics between the three boundary words ($100001$, $100101$, $101001$), or
\item a symmetry action on the two-bit encoding that preserves the folding template while permuting codon indices.
\end{itemize}
Operationally, this can be tested by computing boundary-hit enrichment of the organism's annotated start/stop codon sets under $\mu^\ast$, and comparing against a background codon model (Appendix~\ref{app:stat_tests}).
Sequence-level tests on selected nonstandard-code corpora are reported in Appendix~\ref{app:nonstandard_sequence_tests}.

\subsection{Reproducible spectrum extraction on arbitrary genomes}
\label{sec:reproducible_pipeline}

Appendix~\ref{app:reproducibility} documents a minimal command-line tool that:
\begin{itemize}
\item reads FASTA DNA/RNA;
\item normalizes DNA to RNA by $T\mapsto U$;
\item scans reading frames and (optionally) extracts ORFs using \texttt{AUG} starts and \texttt{UAA/UAG/UGA} stops;
\item outputs per-codon traces $(Z,U)$ as a TSV file (optional) and per-segment summary statistics as JSONL (optional), including boundary-hit counts, $V/\Delta$ histograms, and stop-context uplift windows.
\end{itemize}

This makes the theory falsifiable on any dataset without hidden assumptions.

\paragraph{Stop-context uplift output.}
When \path{--stop-window} is enabled, the CLI reports for each stop class $s\in\{\mathrm{UAA,UAG,UGA}\}$ the before/after uplift window means and histograms (Appendix~\ref{app:stop_context_windows}), directly supporting quantitative tests of prediction~P2.

\subsection{Genome-scale check on human RefSeq mRNA}
\label{sec:refseq_transcriptome_check}

As an initial large-corpus check of predictions~P1--P2 under $\mu^\ast$, we scan the human RefSeq mRNA collection.
For each transcript we select the longest ORF across frames and compute terminal-stop statistics, boundary-hit rates, stop-context uplift windows, and basic $Z$-spectrum fingerprint metrics.
The full tables are recorded in Appendix~\ref{app:refseq_transcriptome}; here we report the headline summary:
On the human RefSeq mRNA corpus (best-ORF per transcript, $\mu^\ast$), we analyzed $n=271181$ transcripts with a detected ORF (out of $n=273423$ records). The terminal stop distribution is $\mathrm{UAA}:78798$, $\mathrm{UAG}:58251$, $\mathrm{UGA}:134132$ (total 271181); the boundary-stop rate is 78798/271181. Across coding tokens (excluding terminal stops), the boundary rate is $\widehat{p}_B=0.0777$. ORF length (codons, excluding stop): mean 573.8, median 407, IQR [161,738].


\paragraph{Terminal stop usage and the boundary-stop rate.}
At window length $m=6$, only $\mathrm{UAA}$ is boundary-aligned (Section~\ref{sec:stop_fine_structure}).
Thus the \emph{boundary-stop rate at termination} in a corpus is exactly the empirical frequency of terminal $\mathrm{UAA}$.
In human RefSeq mRNA, $\mathrm{UGA}$ is the most frequent terminal stop, followed by $\mathrm{UAA}$ and $\mathrm{UAG}$ (Appendix~\ref{app:refseq_transcriptome}).

\paragraph{Stop-context uplift windows.}
Appendix~\ref{app:stop_context_windows} defines the before/after uplift window means, and Appendix~\ref{app:refseq_transcriptome} reports their empirical values for terminal stop codons at window radius $k=10$.
The corresponding pairwise Welch tests are:
Welch tests (two-sided) for stop-context window means (window radius $k=10$) across terminal stops in the RefSeq corpus. Before-window: $p(\mathrm{UAA}\!\neq\!\mathrm{UAG})<10^{-300}$, $p(\mathrm{UAA}\!\neq\!\mathrm{UGA})<10^{-300}$, $p(\mathrm{UAG}\!\neq\!\mathrm{UGA})=0.0052$. After-window: $p(\mathrm{UAA}\!\neq\!\mathrm{UAG})=0.0024$, $p(\mathrm{UAA}\!\neq\!\mathrm{UGA})<10^{-300}$, $p(\mathrm{UAG}\!\neq\!\mathrm{UGA})<10^{-300}$.

These differences are quantitative evidence that terminal-stop neighborhoods carry distinct uplift signatures across stop classes in the human transcriptome.
Effect-size summaries and multi-$k$ sensitivity tables are recorded in Appendix~\ref{app:refseq_transcriptome}.
To separate uplift effects from nucleotide-composition biases (GC and dinucleotide usage), we also report composition-matched controls as defined in Appendix~\ref{app:composition_controls}.

\paragraph{Start-context uplift windows.}
Appendix~\ref{app:start_context_windows} defines the analogous before/after uplift windows around the start codon.
We report start-context window means and multi-$k$ sensitivity in Appendix~\ref{app:refseq_transcriptome}.

\paragraph{Codon-usage bias relative to folding spectra.}
Beyond control sites, the folding map assigns every coding token a pair $(Z,U)$.
In the same RefSeq scan, we compare the observed corpus averages $(\overline{Z},\overline{U})$ to an amino-acid preserving null that draws synonymous codons uniformly within each amino-acid family.
We find a strong deviation under this null:
Codon-usage summary over coding tokens (excluding terminal stops): $\overline{Z}=8.0661$, $\overline{U}=20.9385$. Under an amino-acid preserving null (uniform choice among synonymous codons), $\mathbb{E}[\overline{Z}]=8.0505$ with $\mathrm{sd}=0.000295$ ($z=52.92$, $p<10^{-300}$), and $\mathbb{E}[\overline{U}]=21.5696$ with $\mathrm{sd}=0.000699$ ($z=-903.43$, $p<10^{-300}$).

This indicates that codon usage in human RefSeq mRNA is not neutral with respect to the folding-derived spectra, even after conditioning on amino-acid composition.
An additive decomposition of the deviation across amino-acid families and synonymous codons is recorded in Appendix~\ref{app:refseq_transcriptome}.

\paragraph{\texorpdfstring{$Z$}{Z}-spectrum fingerprint metrics.}
Finally, we report simple per-ORF summary features---boundary-hit rate, entropy of $Z$, and lag-1 autocorrelation of $Z$---as an initial sequence-level fingerprint family:
\begin{tabular}{@{}l@{}}
Z-spectrum fingerprint metrics over best ORFs (excluding terminal stops):\\
boundary-rate mean 0.0857, median 0.0799\\
entropy $H(Z)$ mean 4.1137, median 4.1663\\
lag-1 autocorrelation $\rho(Z_i,Z_{i+1})$ mean -0.0047, median 0.0002 (n=270844).
\end{tabular}

The near-zero lag-1 autocorrelation in particular suggests that the minimal $Z$ process is close to uncorrelated at the codon scale in this dataset, while the entropy values indicate broad usage of the available $Z$ levels.

\subsection{Cross-domain corpus panel replication}
\label{sec:corpus_panel_replication}

To move beyond a single eukaryotic transcriptome and to support systematic replication across domains, we evaluate a cross-domain corpus panel (Eukaryota/Bacteria/Archaea) under the same $\mu^\ast$ pipeline.
The panel summary and per-item tables (including start-context windows and stop-context effect sizes across datasets) are recorded in Appendix~\ref{app:corpus_panel}.
Corpus panel \path{corpus_panel_v2} generated $n=8$ item summaries with $k\in\{3,5,10,20\}$.


\subsection{Human transcript case studies: HBB and INS}
\label{sec:human_case_studies}

As a reproducible sanity check at the sequence level, we apply the pipeline to two human transcripts:
HBB (RefSeq \path{NM_000518.5}) and INS (Ensembl \path{ENST00000381330.5}) \cite{NCBIRefSeqNM000518,EnsemblENST00000381330}.
For each transcript we select the longest ORF across the three frames and record its start/stop positions and folded signatures under $\mu^\ast$.

\begin{center}
\scriptsize
\setlength{\tabcolsep}{4pt}
\renewcommand{\arraystretch}{1.15}
\resizebox{\textwidth}{!}{%
\begin{tabular}{lrrrrrlll}
\toprule
transcript & nt & codons & frame & start (1-based) & stop (1-based) & stop codon & start $(N,w,V,\Delta)$ & stop $(N,w,V,\Delta)$ \\
\midrule
HBB (\path{NM_000518.5}) & 628 & 148 & 2 & 51 & 492 & UAA & $(14,\mathtt{100001},14,0)$ & $(48,\mathtt{100001},14,34)$ \\
INS (\path{ENST00000381330.5}) & 465 & 111 & 2 & 60 & 390 & UAG & $(14,\mathtt{100001},14,0)$ & $(50,\mathtt{001001},16,34)$ \\
\bottomrule

\end{tabular}
}
\end{center}

The HBB transcript terminates with $\mathrm{UAA}$ and exhibits the boundary-aligned $14/48$ start--stop split.
The INS transcript terminates with $\mathrm{UAG}$ and exhibits a non-boundary stop signature, consistent with the stop fine structure in Section~\ref{sec:stop_fine_structure}.
In both cases, the reported start/stop tuples match the corresponding entries in the control codon table (Section~\ref{sec:control_codon_table}).


